\section{Discussion and Limitations}\label{sec:discussion}

\paragraph{What the safety results protect.}
The floor-cap theorem and the credit-capacity results concern a balance sheet in units of the reserve asset. They give no protection against a fall in the value of that asset, a reserve token that stops being transferable, or reserve investments that cannot be redeemed at once. An earning reserve changes the opportunity cost of lending; a temporarily illiquid reserve fails the immediate cash requirement for as long as it is unavailable.

\paragraph{The economic cost of additional credit.}
The floor is a boundary imposed by a combination of promises: proportional collateral release, independent repayment and no additional committed reserve. \Cref{sec:credit} shows exactly how relaxing one of those promises permits more credit. A buffer commits cash that would otherwise be withdrawable. A curve-valued lock restricts release and splitting, and needs stronger guards when the curve changes. These are different allocations of cash and withdrawal rights. Greater nominal borrowing capacity does not establish a welfare gain, a positive borrower residual, or demand for the resulting contract. The original floor cap remains attractive for its simple local check, permissive release and robustness to every funded reshape that preserves the floor. In stepped designs, borrowers whose locks fit within the floor capacity gain no extra credit from the nonlinear cap.

\paragraph{Where a higher floor comes from.}
Without accrual, reserve for a floor raise comes from a buyer's premium, a trader's fee, an origination fee or an external contribution. A fee-funded raise transfers resources from fee payers to those holding tokens at the time. Capped accrual can instead fund a raise through increased borrower obligations (\cref{sec:carry}), reducing borrowers' residual value. These accounting operations do not create additional financial resources. A platform can nevertheless provide a valuable service, and liquidity can matter when users face financing frictions. Such benefits need to be modeled and funded. \Cref{prop:deposit-benchmark} shows why borrowing against a redeemable service token must be compared with a refundable deposit offering the same service.

\paragraph{How far the aggregate results reach.}
\Cref{thm:eq} determines borrowing plans and aggregate quantities at the floor with given arrivals and horizons. \Cref{cor:noentry} shows that cash-only agents would not supply positive arrivals when growth is below their discount rate and the cap is below the net floor. Above the floor, exit order changes redemption proceeds and decisions become interdependent. A model of participation and leveraged exposure under stated beliefs, rewards and any service value is still needed. The case $g\ge\rho$, in which the floor outgrows the discount rate on later buyers' fees, also remains outside the analysis. Adding an arbitrary positive benefit would make an entry inequality hold without explaining demand.

\paragraph{What the evidence establishes.}
The Lean development covers the real-valued floor-cap invariant. It does not cover the new credit-capacity results, the integer refinement or a production state machine. The original Solidity implementation is tested, not formally verified or audited; the proposed alternative contracts are not implemented there. Exact-rational checks reproduce the new examples and compare independent finite algorithms, while the quantified claims rest on their analytic proofs. The empirical study concerns one protocol, six pools and seven months, with no floor movement in five pools. Fee allocations and claim events document separate parts of the reward mechanism; they do not measure account earnings, borrower rebates or returns. Neither the observed use of leverage nor the repayment patterns identify a causal response to the contract's terms.

\paragraph{Disclosure.} The author is a co-founder of House Markets, which develops token issuance and credit mechanisms of the kind analyzed here, and is a co-author of \vone, on which this paper builds (\cref{app:errata}); the authors of that paper were affiliated with Floors Finance and House Markets. The author holds no tokens of Baseline or Nirvana, the two deployed systems discussed. The data, the scripts and every number of \cref{sec:empirics} are in the supplementary material.

\paragraph{Use of AI tools.} During the preparation of this work, the author used AI assistants (Anthropic Claude and OpenAI Codex) to draft and edit text, write and test code, construct and check proofs, including the Lean~4 formalization, and review the manuscript. The author reviewed, edited and verified all output, and takes full responsibility for the content of this paper.

\paragraph{Further tests that would change the conclusions.}
Three directions follow directly from the results. First, an integer implementation and machine-checked refinement could establish whether the new guards remain tractable and safe with storage, rounding and token transfers. Second, an endogenous participation model should compare the mechanism with deposits and outside finance while paying for service costs; its prediction should distinguish leveraged exposure from demand for liquidity. Third, a broader event study with complete account inventories and reward accounting could estimate realized outcomes and test responses to identifiable changes in terms. The present data do not supply that identification. The theoretical capacity gain should also be measured against actual position sizes and floor capacity before it is claimed to be economically material.

\section{Conclusion}

Lending the redemption reserve is compatible with redeeming every unlocked token because the collateral is removed from circulation. The floor-cap mechanism makes that argument precise through custody, a per-loan floor bound and adequate cash plus debt, maintained by a complete set of operation guards. The preservation of those conditions is machine-checked for the real-valued model.

The larger design question is which future collateral withdrawals a loan promises. Current aggregate coverage does not answer it. With proportional release and no committed buffer, every loan must respect the floor value of its collateral. A sorted residual-debt envelope gives the exact buffer for exceeding that bound and the maximum aggregate credit for a fixed reserve. A curve-valued loan cap can use more of the collateral's reserve value, but changes partial release, splitting and permissible repricing. Additional credit therefore has a precise contractual or capital cost.

These safety results leave the participation question open. Under the floor-cap assumptions, long-horizon holders gain by borrowing and deferring repayment, yet the stationary floor economy cannot attract cash-only entrants when growth is below their discount rate and the loan cap is below the net floor. Service value can cover the capital cost, but an equivalent refundable deposit is a relevant benchmark. The Baseline sample documents substantial leveraged purchasing and fee allocations to staking. It does not establish the beliefs or net benefits behind those choices. The paper supplies conditions for safe obligations and sharper tests for a demand theory; it does not infer economic success from accounting safety.
